\section{Estimation du projet}

% ----------------------------------------------------------------
% 📌 À quoi sert cette section ?
% ----------------------------------------------------------------
% Cette partie permet d’évaluer de manière réaliste le volume de travail nécessaire pour réaliser le projet.
% Elle montre que vous avez réfléchi à la charge de travail, aux ressources humaines, et au temps disponible.
% Cette estimation permet aussi au jury de vérifier que le projet correspond bien aux attentes du 4LABO.

% ----------------------------------------------------------------
% 🧭 Comment la remplir ?
% ----------------------------------------------------------------
% ➤ Rappel : 1 jour-homme (JH) = 2 points 4LABO
% ➤ Additionnez les estimations de toutes les fonctionnalités prévues
% ➤ Mentionnez le nombre de porteurs pour répartir l’effort
% ➤ Soyez cohérents avec les User Stories et la complexité réelle du projet
% ➤ Utilisez un tableau si besoin pour justifier le total

% ----------------------------------------------------------------
% 🧾 Exemple rédigé :
% ----------------------------------------------------------------
% Le projet est porté par deux étudiants.
% En se basant sur les User Stories détaillées ci-dessus, nous estimons le travail total à 12 JH,
% soit 6 JH par étudiant.
%
% Cela correspond à une charge totale de \textbf{12 JH = 24 points 4LABO},
% donc \textbf{12 points 4LABO par étudiant}.

% ----------------------------------------------------------------
% 📊 Exemple de tableau récapitulatif :
% ----------------------------------------------------------------
% \begin{table}[H]
% \centering
% \begin{tabular}{|l|c|}
% \hline
% Élément estimé & JH \\
% \hline
% Fonctionnalités principales & 6 \\
% Documentation utilisateur & 2 \\
% Déploiement et tests & 2 \\
% Vidéo de démonstration & 2 \\
% \hline
% \textbf{Total estimé} & \textbf{12} \\
% \hline
% \end{tabular}
% \caption{Estimation globale du projet}
% \end{table}

% ----------------------------------------------------------------
% ✍️ À vous de jouer : complétez ci-dessous
% ----------------------------------------------------------------

% (Suggestion de départ)
Nous estimons que le projet nécessite environ :

\begin{itemize}
  \item [À compléter...]
\end{itemize}
